\subsection{Essential image of a measurable function}

In this issue, we denote by X a locally compact topological space, and by \(\mu\) a positive measure on X.

DEFINITION 1. --- Let Y be a topological space. For every measurable function g from X to Y, the \(\mu\)-essential image of \(g\) is the subset of the \(y \in Y\) such that, for every open neighborhood \(U\) of \(y\), the set \(g^{-1}(\mathrm{U})\) is not locally \(\mu\)-negligible in X (INT, IV, p. 172, § 5, no. 2, def. 3).

Lemma 1. --- Let \(g\) be a \(\mu\)-measurable function from X into a topological space Y and let \(S\) be its \(\mu\)-essential image. The set of elements \(x \in X\) such that \(g(x)\) does not belong to S is locally \(\mu\)-negligible in X.

Let \(Z = g^{-1}(Y - S)\) be the set in question.

Suppose first that X is compact and g is continuous. In this case, the pushforward measure \(g(\mu)\) is defined (INT, V, p. 69, § 6, no. 1, def. 1) since \(\mu\) is a bounded measure. It follows from the definitions that the \(\mu\)-essential image of \(g\) is the support of the measure \(g(\mu)\) (cf. INT, V, p. 70, § 6, no. 2, cor. 2), whence \(\mu(Z) = g(\mu)(Y - S) = 0\) (INT, IV, p. 118, § 2, no. 2, prop. 5).

Now consider the general case. Let C be a compact subset of X, and let \(\varepsilon > 0\) be a real number. Since \(g\) is measurable, there exists a compact subset \(C_1 \subset C\) such that \(\mu(C - C_1) \leqslant \varepsilon\) and such that \(g|C_1\) is continuous (INT, IV, p. 169, § 5, no. 1, prop. 1). We then have

\[
\mu (\mathrm{Z} \cap \mathrm{C}) \leqslant \mu (\mathrm{C} - \mathrm{C} _ {1}) + \mu (\mathrm{Z} \cap \mathrm{C} _ {1}) \leqslant \varepsilon + \mu (\mathrm{Z} \cap \mathrm{C} _ {1}).
\]

Let \(\mu_{1}\) be the measure induced by \(\mu\) on \(C_{1}\) (INT, IV, p. 186, § 5, no. 7, def. 4). Let \(S_{1}\) be the \(\mu_{1}\)-essential image of \(g|C_{1}\) . We have \(Z\cap C_{1}\subset(g|C_{1})^{-1}(Y-S_{1})\) . By the first case, the set \(Z\cap C_{1}\) is therefore \(\mu_{1}\) -negligible, and consequently \(\mu\) -negligible (INT, IV, p. 187, § 5, no. 7, lemma 2, (i)). The above inequality becomes \(\mu(Z\cap C)\leqslant\varepsilon\) ; since \(\varepsilon\) and C are arbitrary, the set Z is locally \(\mu\) -negligible (INT, IV, p. 172, § 5, no. 2, prop. 5).

Lemma 2. --- Let \(g\) be a continuous function from X into \(\mathbf{C}\). The \(\mu\)-essential image of \(g\) is the closure of the image of the support of \(\mu\) under \(g\).

Let \(Y\) be the support of \(\mu\). If \(z \in \mathbf{C}\) is not adherent to \(g(Y)\), then there exists an open neighborhood \(U\) of \(z\) such that the open set \(g^{-1}(U)\) does not meet Y and is therefore locally \(\mu\)-negligible; this means that \(z\) does not belong to the \(\mu\)-essential image of \(g\) .

Conversely, if \(z \in \mathbf{C}\) is adherent to \(g(Y)\), then for every open neighborhood \(U\) of z, the set \(g^{-1}(U)\) is an open set in X which meets Y. It is therefore not \(\mu\)-negligible, and since it is open, it is not locally \(\mu\)-negligible (INT, IV, p. 172, § 5, n°2, cor. 2). Hence z belongs to the \(\mu\)-essential image of g.

Lemma 3. --- Let \(g\) be a continuous function from X into \(\mathbf{C}\) and let S be its \(\mu\)-essential image. Suppose that S is nonempty and that \(0 \notin S\), and denote by \(\delta\) the distance from 0 to S. Then \(\delta > 0\).

Set \(h(x) = 1 / g(x)\) if \(g(x) \neq 0\) and \(h(x) = 0\) otherwise. The function \(h\) belongs to \(\mathcal{L}^{\infty}(\mathrm{X},\mu)\). Let \(\widetilde{h}\) be the class of \(h\) in \(\mathrm{L}^{\infty}(\mathrm{X},\mu)\) . We then have the formula \(\delta^{-1} = \| \widetilde{h}\|_{\infty}\) .

Let U be an open neighborhood of 0 such that the open set \(Z = g^{-1}(U)\) is locally \(\mu\) -negligible, hence negligible (INT, IV, p. 172, § 5, no. 2, cor. 2). Let Y be the support of \(\mu\) ; we have \(Y \subset X - Z\) . The restriction of the function \(h\) to \(X - Z\) is continuous and bounded, hence \(h \in \mathcal{L}^{\infty}(X, \mu)\) and the norm of \(\widetilde{h}\) in \(\mathcal{L}^{\infty}(X, \mu)\) is equal to the norm of its restriction to \(X - Z\) . Moreover, for every \(\alpha \in R_{+}\) the set of \(x \in X - Z\) such that \(|h(x)| > \alpha\) is open in \(X - Z\) ; it is therefore locally \(\mu\) -negligible if and only if it does not meet Y (INT, IV, loc. cit. and INT, III, p. 66, § 3, no. 2, def. 1). Consequently, we have

\[
\| \widetilde {h} \| _ {\infty} = \sup _ {x \in \mathrm{Y}} \frac {1}{| g (x) |},
\]

whence

\[
\frac {1}{\| \widetilde {h} \| _ {\infty}} = \inf _ {x \in \mathrm{Y}} | g (x) | = \inf _ {\lambda \in g (\mathrm{Y})} | \lambda | = \inf _ {\lambda \in \overline {{g (\mathrm{Y})}}} | \lambda |
\]

which is equal to δ by the preceding lemma.

